% GENERATED FILE. Edit canonical lessons, then run npm run generate:books.
% canonical-source-hash: fnv1a64:af0b442f69ac05f6
% canonical-lessons: UR-W35-jim, UR-W35-bari-he, UR-W35-khe, UR-W35-che

\chapter{The Hook Family}
\label{ch:ur-letters-hook-family}
\hlchaptermodality{35}

\begin{chapteropening}
\textbf{By the end of this chapter:} \emph{I can write jim, bari he, khe and che, and tell the four apart by their dots.}

Complete the last lesson of chapter 35: i can write jim, bari he, khe and che, and tell the four apart by their dots.
\end{chapteropening}

\section[jīm]{\ur{ج} (jīm) — the letter that opens \ur{جی}}

\label{lesson:UR-W35-jim}



\begin{quote}
Before the new one: write \ur{ت} and \ur{ٹ} once each, and name them.

Now say \textbf{\ur{جی} \ur{ہاں}} (\emph{jī hā̃}), \textquotedblleft{}yes\textquotedblright{}. You have read it for a long time. This lesson takes one letter out of it and puts it in your hand.
\end{quote}

\subsection*{Script you'll notice: \ur{ج}}
\textbf{\ur{ج}} — \emph{jīm}, the sound \emph{j}.

It opens \textbf{\ur{جی}} (\emph{jī}) in \textbf{\ur{جی} \ur{ہاں}}, \textquotedblleft{}yes\textquotedblright{}, from the first chapter. A hooked head, a deep bowl, and \textbf{one dot below}.

\subsection*{Writing: \ur{ج}}
\hlblockfigure{\detokenize{figures/UR-W35-jim-filmstrip.pdf}}{How \ur{ج} is written, stroke by stroke}

\begin{itemize}
\raggedright
  \item \textbf{1.} start here: place the dot below
  \item \textbf{2.} lift, then sweep left through the pointed hooked head
  \item \textbf{3.} without lifting, continue down and around the bowl — and only now lift
\end{itemize}

\textbf{Pen lifts: 1.}

\begin{quote}
Stroke order is one attested teaching order; hands and teachers vary. Source: Zer o Zabar, Te, mim, jim, che, and more diacritics, independent \ur{ج} handwriting animation and flat-head insight (Northwestern University Libraries Open Textbook, 2026).
\end{quote}

\subsection*{Your turn}
\begin{itemize}
\raggedright
  \item \emph{Look:} at \textbf{\ur{جی} \ur{ہاں}} and point at \ur{ج} inside it
  \item \emph{Trace it:} \ur{ج} three times, saying \emph{j} as you finish each one
  \item \emph{Write it:} \ur{ٹ}, then \ur{ج}, from memory
  \item \emph{Read it:} \textbf{\ur{جی} \ur{ہاں}} aloud once more
\end{itemize}

\begin{tcolorbox}[breakable,colback=teal!4,colframe=teal!35!black,title={Before you move on}]
Which letter is this — \ur{ج}? (\textbf{jīm}, \emph{j}.) Which word did it come from? (\textbf{\ur{جی} \ur{ہاں}}, \emph{jī hā̃}, \textquotedblleft{}yes\textquotedblright{}.)
\end{tcolorbox}
\section[baṛī he]{\ur{ح} (baṛī he) — the letter that opens \ur{حافظ}}

\label{lesson:UR-W35-bari-he}



\begin{quote}
Before the new one: write \ur{ٹ} and \ur{ج} once each, and name them.

Now say \textbf{\ur{حافظ}} (\emph{hāfiz}), \textquotedblleft{}guardian\textquotedblright{}. You have read it for a long time. This lesson takes one letter out of it and puts it in your hand.
\end{quote}

\subsection*{Script you'll notice: \ur{ح}}
\textbf{\ur{ح}} — \emph{baṛī he}, \textquotedblleft{}big \emph{he}\textquotedblright{}, the sound \emph{h}.

It opens \textbf{\ur{حافظ}} (\emph{hāfiz}), the second half of \textbf{\ur{خدا} \ur{حافظ}}. Urdu already has an \emph{h}, \textbf{\ur{ہ}}, which you write; this one came in with Arabic words and sounds the same. It is the hook-and-bowl skeleton of \textbf{\ur{ج}}, with \textbf{no dots}.

\subsection*{Writing: \ur{ح}}
\hlblockfigure{\detokenize{figures/UR-W35-bari-he-filmstrip.pdf}}{How \ur{ح} is written, stroke by stroke}

\begin{itemize}
\raggedright
  \item \textbf{1.} start here: sweep left through the pointed hooked head
  \item \textbf{2.} continue down and around the deep bowl without lifting — and only now lift
\end{itemize}

\textbf{Pen lifts: 0.} The pen never leaves the paper.

\begin{quote}
Stroke order is one attested teaching order; hands and teachers vary. Source: Zer o Zabar, Sīn, shīn, baṛī he, nūn, and nūn-e ġhunna, independent \ur{ح} handwriting animation and Baṛī he instructions (Northwestern University Libraries Open Textbook, 2026).
\end{quote}

\subsection*{Your turn}
\begin{itemize}
\raggedright
  \item \emph{Look:} at \textbf{\ur{حافظ}} and point at \ur{ح} inside it
  \item \emph{Trace it:} \ur{ح} three times, saying \emph{h} as you finish each one
  \item \emph{Write it:} \ur{ج}, then \ur{ح}, from memory
  \item \emph{Read it:} \textbf{\ur{حافظ}} aloud once more
\end{itemize}

\begin{tcolorbox}[breakable,colback=teal!4,colframe=teal!35!black,title={Before you move on}]
Which letter is this — \ur{ح}? (\textbf{baṛī he}, \emph{h}.) Which word did it come from? (\textbf{\ur{حافظ}}, \emph{hāfiz}, \textquotedblleft{}guardian\textquotedblright{}.)
\end{tcolorbox}
\section[khe]{\ur{خ} (khe) — the letter that opens \ur{خدا}}

\label{lesson:UR-W35-khe}



\begin{quote}
Before the new one: write \ur{ج} and \ur{ح} once each, and name them.

Now say \textbf{\ur{خدا}} (\emph{khudā}), \textquotedblleft{}God\textquotedblright{}. You have read it for a long time. This lesson takes one letter out of it and puts it in your hand.
\end{quote}

\subsection*{Script you'll notice: \ur{خ}}
\textbf{\ur{خ}} — \emph{khe}, a \emph{kh} scraped at the back of the throat.

It opens \textbf{\ur{خدا}} (\emph{khudā}), the first half of \textbf{\ur{خدا} \ur{حافظ}}. The same skeleton as \textbf{\ur{ح}}, with \textbf{one dot above}.

\subsection*{Writing: \ur{خ}}
\hlblockfigure{\detokenize{figures/UR-W35-khe-filmstrip.pdf}}{How \ur{خ} is written, stroke by stroke}

\begin{itemize}
\raggedright
  \item \textbf{1.} start here: draw the short upper head from left to right
  \item \textbf{2.} without lifting, continue down and around the deep bowl
  \item \textbf{3.} lift once, then place the dot above — and only now lift
\end{itemize}

\textbf{Pen lifts: 1.}

\begin{quote}
Stroke order is one attested teaching order; hands and teachers vary. Source: Zer o Zabar, Khe, ze, zāl, swād, and zwād, independent \ur{خ} handwriting animation and Khe instructions (Northwestern University Libraries Open Textbook, 2026).
\end{quote}

\subsection*{Your turn}
\begin{itemize}
\raggedright
  \item \emph{Look:} at \textbf{\ur{خدا}} and point at \ur{خ} inside it
  \item \emph{Trace it:} \ur{خ} three times, saying \emph{kh} as you finish each one
  \item \emph{Write it:} \ur{ح}, then \ur{خ}, from memory
  \item \emph{Read it:} \textbf{\ur{خدا}} aloud once more
\end{itemize}

\begin{tcolorbox}[breakable,colback=teal!4,colframe=teal!35!black,title={Before you move on}]
Which letter is this — \ur{خ}? (\textbf{khe}, \emph{kh}.) Which word did it come from? (\textbf{\ur{خدا}}, \emph{khudā}, \textquotedblleft{}God\textquotedblright{}.)
\end{tcolorbox}
\section[che]{\ur{چ} (che) — the letter inside \ur{سوچنا}}

\label{lesson:UR-W35-che}



\begin{quote}
Before the new one: write \ur{ح} and \ur{خ} once each, and name them.

Now say \textbf{\ur{سوچنا}} (\emph{sochnā}), \textquotedblleft{}to think\textquotedblright{}. You have read it for a long time. This lesson takes one letter out of it and puts it in your hand.
\end{quote}

\subsection*{Script you'll notice: \ur{چ}}
\textbf{\ur{چ}} — \emph{che}, the sound \emph{ch}.

It is inside \textbf{\ur{سوچنا}} (\emph{sochnā}), \textquotedblleft{}to think\textquotedblright{}. The same skeleton once more, with \textbf{three dots below}: one letter, four sounds, and the dots decide which.

\subsection*{Writing: \ur{چ}}
\hlblockfigure{\detokenize{figures/UR-W35-che-filmstrip.pdf}}{How \ur{چ} is written, stroke by stroke}

\begin{itemize}
\raggedright
  \item \textbf{1.} start here: sweep left through the pointed hooked head
  \item \textbf{2.} continue down and around the bowl without lifting
  \item \textbf{3.} after one lift, place the lower-left dot
  \item \textbf{4.} after another lift, place the lower-right dot
  \item \textbf{5.} after a third lift, place the lower-center dot — and only now lift
\end{itemize}

\textbf{Pen lifts: 3.}

\begin{quote}
Stroke order is one attested teaching order; hands and teachers vary. Source: Zer o Zabar, Te, mīm, jīm, che, and more diacritics, independent \ur{چ} calligraphic and handwriting animations and Che instructions (Northwestern University Libraries Open Textbook, 2026).
\end{quote}

\subsection*{Your turn}
\begin{itemize}
\raggedright
  \item \emph{Look:} at \textbf{\ur{سوچنا}} and point at \ur{چ} inside it
  \item \emph{Trace it:} \ur{چ} three times, saying \emph{ch} as you finish each one
  \item \emph{Write it:} \ur{خ}, then \ur{چ}, from memory
  \item \emph{Read it:} \textbf{\ur{سوچنا}} aloud once more
\end{itemize}

\begin{tcolorbox}[breakable,colback=teal!4,colframe=teal!35!black,title={Before you move on}]
Which letter is this — \ur{چ}? (\textbf{che}, \emph{ch}.) Which word did it come from? (\textbf{\ur{سوچنا}}, \emph{sochnā}, \textquotedblleft{}to think\textquotedblright{}.)
\end{tcolorbox}
